\subsubsection{TPTP.}
Two subsets of the TPTP problem library~\cite{Sut17} are used. The first is a
737-problem provable subset (FOF/THM), restricted to problems with 0--50 input
formulae, 0--150 atoms, and no equality or arithmetic. The second is a
204-problem counter-satisfiable subset (FOF/CSA) drawn from the same
complexity bounds, in which each problem is known to have no proof because
its negation is satisfiable. The two subsets together test both sides
of the prover against problems with
known status across diverse mathematical domains.

\subsubsection{FOLIO.}
The FOLIO dataset~\cite{han2024folio} provides 280 human-authored
natural-language reasoning problems formalised as FOF, testing whether the
prover's behaviour generalises beyond the controlled TPTP filter.

\subsubsection{Synthetic.}
202 synthetic FOF problems were generated programmatically across four tiers,
\emph{easy} (32), \emph{medium} (60), \emph{hard} (60), \emph{expert}
(50), with proof depth $n\in[2,15]$ and distractor count $d\in[0,14]$
scaling jointly with tier. Four problem families are used.
\textbf{Implication chains}: a conjecture reachable from a ground fact via
$n$ universally quantified steps, with $d$ structurally identical
distractor axioms. \textbf{Syllogisms}: a fixed class hierarchy with a
ground instance for one individual. \textbf{Transitivity chains}: a single
binary transitivity axiom plus ground edge facts, requiring repeated reuse
of the same axiom across different instantiations.
\textbf{Quantifier alternation}: single-formula conjectures testing
classical $\forall$/$\exists$ theorems and non-theorems. Unprovable
variants negate the conjecture (e.g., remove a chain link, change the
syllogism individual, or ask for an unreachable node).
